\chapter{How Much Do You Want?}

\section{From Building to Stopping}

The preceding chapters construct. They begin with a mark, carry it through receipt,
count, residue, observer, value, scale, support, energy, boundary, radiation,
completion, and orientation, and they keep the single invariant visible through
each change of frame. That path can always invite one more reading. A reader can
ask for a plainer image, a finite device, a theorem statement, a construction,
an analytic reading, or a metaphysical one. If the book answered each further
request in the main line, it would never stop.

This chapter adds no new object to the construction. It gives the construction a
stop rule for its own exposition. The rule asks how much of the explanation
belongs in the main text for a given observer frame and a requested depth. It
does not decide how every theory should be explained. It does not rank readers.
It only says when this finite construction owes the next layer and when the next
layer belongs outside the main line.

The question is therefore literal: how much do you want? The answer is not an
unbounded invitation. The answer is a finite projection, followed by a single
on/off decision about whether the next depth changes what the main text owes.
The rule does not shame a deeper request, and it does not flatten a difficult
construction into a slogan. It separates two acts that the book has to keep
apart: giving enough for the present frame, and preserving a place where the
next frame can be read without forcing itself into the present one.

\section{The Observer Frame and the Depth Ladder}

The invariant does not change because a different reader asks for a different
depth. What changes is the observer frame: the kind of reading requested, the
amount of construction exposed, and the boundary between the main text and the
appendix. A parable can carry the same invariant as a finite device. A theorem
statement can carry the same invariant as a longer construction. An analytic
or metaphysical reading can extend the account without replacing the thing
being carried.

The book therefore uses a depth ladder. The ladder is not a hierarchy of truth.
It is a way to pitch the same construction at a chosen explanatory depth:

\[
\begin{array}{c|l}
d & \text{reading depth} \\
\hline
0 & \text{parable} \\
1 & \text{finite machine} \\
2 & \text{theorem statement} \\
3 & \text{construction} \\
4 & \text{analytic / physical reading} \\
5+ & \text{metaphysical / topological reading}
\end{array}
\]

Depth 0 gives the reader the image. Depth 1 gives the finite device. Depth 2
states the theorem-shaped claim. Depth 3 shows how the construction is made.
Depth 4 asks how the construction faces analytic or physical language. Depth
5 and above ask how the construction reads as metaphysics or topology. These
depths do not compete. They are projections of the same finite account.

Most of the main line lives at depths 0 through 2. A deeper passage earns its
place only when it changes what the reader must know in order to keep the
construction honest. Otherwise the deeper reading can exist, but it does not
belong in the main line.

The observer frame matters because a reader does not merely receive more or
less information. A reader receives a kind of account. A general technical
reader may need the parable and the finite device. A mathematical reader may
need the theorem statement and enough construction to see why the statement is
not decorative. A reader following the analytic or physical face may need the
limits on that reading stated without being asked to carry the metaphysical
account at the same time. A philosophical reader may ask for the deeper
topological grammar. The same invariant can survive each request because the
depth changes the presentation, not the object carried through the presentation.

\section{Projection and the Appendix Boundary}

Fix an object of exposition \(E\), and let \(d\) be the depth requested by the
observer frame. The main text projects the explanation to that depth. Material
above that depth begins in the appendix:

\[
\begin{aligned}
\Pi_{\le d}(E) &=
  \text{the explanation of }E\text{ through depth }d,\\
\operatorname{Appendix}(E,d) &= \Pi_{\ge d+1}(E).
\end{aligned}
\]

The projection does not discard the higher reading. It gives the main text a
boundary. Everything through depth \(d\) belongs to the selected explanation.
Everything beginning at \(d+1\) waits outside the main line unless it earns
entry by changing what the next frame owes.

The appendix boundary is therefore not a wastebasket. It is the place where
owed-but-not-present material can remain available without distorting the
current frame. A construction can keep its proof-shaped detail there when the
main text is operating as theorem statement. An analytic aside can wait there
when the main text is carrying only the finite device. A metaphysical reading
can wait there when it would turn a local distinction into a larger meditation.
The point is not to hide the deeper account. The point is to keep each account
from impersonating the depth beside it.

This boundary keeps the book finite. Without it, every theorem statement would
drag its full construction behind it, every construction would drag its analytic
reading, and every analytic reading would drag a metaphysical commentary. The
main text would confuse depth with completeness. The stop rule keeps those
roles apart. A reader may ask for more, but the book does not silently fold
every higher reading into the present one.

The appendix boundary also protects the reader who asks for less. A reader at
depth 1 does not need every depth 4 consequence in the main line. A reader at
depth 3 does not need every depth 5 interpretation before the construction can
stand. The projection says what is being given now; the appendix says where the
rest begins.

This also explains why the rule belongs at the end rather than at the
beginning. At the beginning, a stop rule would sound like a promise made before
the construction had earned its terms. At the end, the book has already shown
the reader marks, counts, residues, observers, boundaries, and orientation. The
reader has seen enough of the apparatus to understand why a depth selector is
not an excuse to avoid work. It is a final boundary drawn by a construction that
has learned how to draw boundaries.

\section{The On/Off Stop Rule}

After projecting to depth \(d\), the book asks one question about depth \(d+1\):
does the next level add something the main text owes? There are four cases.
Only three turn the signal on:

\[
\begin{array}{c|c|c}
\text{next level adds} & \text{signal} & \text{main-text action} \\
\hline
\varnothing & \mathsf{off} & \text{stop at depth }d \\
\text{measurable prediction} & \mathsf{on} & \text{open the next frame} \\
\text{formal theorem} & \mathsf{on} & \text{open the next frame} \\
\text{necessary distinction} & \mathsf{on} & \text{open the next frame}
\end{array}
\]

If the next level adds no measurable prediction, no formal theorem, and no
necessary distinction, the signal is off. The main text stops at depth \(d\),
and the deeper material belongs in the appendix or in a later reading. If the
next level adds any one of those three things, the signal is on. The next frame
has earned a place because omitting it would hide something the construction
needs.

The three on-cases are deliberately few. A measurable prediction changes what a
reader can check. A formal theorem changes what the construction has actually
secured. A necessary distinction changes what the reader must keep apart in
order not to confuse the claim. Preference, curiosity, ornament, and associative
pressure do not turn the signal on by themselves.

The rule is strict because the ending is a dangerous place. An ending tempts the
book to gather every unspent association and call the gathering completion. The
on/off rule refuses that temptation. A next paragraph is not owed because it
would sound satisfying. A next section is not owed because it could repeat the
whole arc at a higher pitch. A next depth is owed only when the construction
would become less honest without it.

This is the same kind of on/off reading the book has used throughout: a finite
boundary reads whether a mark is carried, whether a refusal matters, whether a
needle has selected one side rather than another. Here the needle points at the
book's own depth. More explanation is owed only when the next depth changes the
finite account in one of the three admitted ways.

This makes the stop rule recursive without making it endless. If the next frame
is opened, the same question can be asked again at the next boundary. Does the
following depth add a measurable prediction, a formal theorem, or a necessary
distinction? If yes, the signal is on. If no, the signal is off. The recursion
halts because the book never asks for a highest possible depth. It asks for the
next owed depth, and then it applies the same finite decision again.

\section{The Carrier Algebra and the Close}

The stop rule lives in a two-valued carrier. Let \(a\) and \(b\) be on/off
signals for admitted changes at the next depth. The carrier has negation, meet,
join, and the De Morgan laws:

\[
\neg(\mathsf{on})=\mathsf{off},\quad
\neg(\mathsf{off})=\mathsf{on},\quad
a\wedge b,\quad a\vee b,
\]
\[
\neg(a\wedge b)=\neg a\vee \neg b,\qquad
\neg(a\vee b)=\neg a\wedge \neg b.
\]

The algebra is small because the decision is small. It does not need a continuum
of explanatory pressure. It asks whether the next frame is owed or not owed.
Several reasons can meet or join, and their negations obey the same finite
logic as the other on/off readings in the book.

The invariant survives this projection. For every admitted depth \(d\), the
projection changes the amount of explanation, not the invariant carried by the
explanation:

\[
I(\Pi_{\le d}(E)) = I(E)
\qquad\text{for every admitted depth }d.
\]

This line is project-native. It does not assert a new physical theorem, a new
continuum result, or a universal law of exposition. It says that within this
finite construction, the same invariant remains legible when the explanation is
projected to an admitted depth.

The line also prevents a common mistake. A shorter account is not a different
invariant merely because it hides deeper detail. A longer account is not a
truer invariant merely because it exposes more detail. The invariant is what
the projection carries. The depth is how much of the carrying the reader has
asked to see.

The stop rule also applies to itself. The rule belongs in the main text because
it supplies a necessary distinction: the distinction between what this
construction owes at a requested depth and what can move to the appendix. Once
the rule is stated, another layer would have to add a measurable prediction, a
formal theorem, or a necessary distinction. The rule gives no permission to add
more merely because more can be said. Nothing above this rule adds a measurable
prediction, a formal theorem, or a necessary distinction, so the signal is off,
and the book stops.
